\documentclass[a4paper,english]{lipics-v2021}
\usepackage{mysty}
\usepackage{mydraft}

\title{Branching Higher-Dimensional Automata}

\author{Uli Fahrenberg}{LMF, Paris-Saclay University, Gif-sur-Yvette, France}{}{}{}
\author{Georg Struth}{Sheffield}{}{}{}
\author{Krzysztof Ziemiański}{University of Warsaw, Poland}{}{}{}

% \keywords{%
%   higher-dimensional automaton
% }

% \begin{CCSXML}
%   <ccs2012>
%   <concept>
%   <concept_id>10003752.10003766.10003773</concept_id>
%   <concept_desc>Theory of computation~Automata extensions</concept_desc>
%   <concept_significance>500</concept_significance>
%   </concept>
%   </ccs2012>
% \end{CCSXML}

% \ccsdesc[500]{Theory of computation~Automata extensions}

\Copyright{U.~Fahrenberg etc.}
\ArticleNo{\mbox{}}
\authorrunning{Authors Running (Or Not)}
% \relatedversion{}
% \relatedversiondetails{Full Version}{https://arxiv.org/abs/UPDATE}

\begin{document}

\maketitle

\section{Introduction}

\section{ST-Pomsets}

A \emph{pomset} (with event order) $(P, {<}, {\evord}, \lambda)$ on an alphabet $\Sigma$
consists of a finite set $P$, a strict partial order $<$ on $P$,
an irreflexive and acyclic relation $\evord$ on $P$, and a labeling function $\lambda: P\to \Sigma$,
such that ${<}\cup {\evord}$ is a total relation.

A \emph{pomset with interfaces} (\emph{ipomset}) $(P, {<}, {\evord}, \lambda, S, T)$
is a pomset together with subsets of \emph{interfaces}
$S\subseteq P_{\min}$ and $T\subseteq P_{\max}$ of $<$-minimal, resp.\ $<$-maximal elements.

% We will often use indices $<_P$, $S_P$, etc.\ for clarity.

An ipomset $(P, {<}, {\evord}, \lambda, S, T)$ is \emph{discrete} if ${<}=\emptyset$.
It is a \emph{starter} if, additionally, $T=P$,
and a \emph{terminator} if $S=P$.
We write discrete ipomsets as vectors with dots for interfaces and the event order going downwards; for example,
$\loset{\ibullet a\pibullet \\ \pibullet b\pibullet \\ \ibullet a\ibullet}$ denotes
$(\{e_1, e_2, e_3\}, \emptyset, (e_1\evord e_2\evord e_3), (e_1\mapsto a, e_2\mapsto b, e_3\mapsto a), \{e_1, e_3\}, \{e_3\})$.

Let $\Sigma$ be a fixed alphabet and denote by $\Sigma^\ST$ the (infinite) set of starters and terminators labeled by $\Sigma$.

For elements $x$, $y$ of a pomset $P$ we write $x\ll y$ if $x<y$ and there is no $z$ such that $x<z<y$.

\ufin{Blabla isomorphism blabla quotient blabla equality is isomorphism}

\begin{definition}
  An \emph{ST-pomset} is a pomset $(P, {<}, {\evord}, \lambda)$ on $\Sigma^\ST$ such that for all $x\in P\setminus P_{\min}$,
  $S_{\lambda(x)}=\bigotimes_{y\ll x} T_{\lambda(y)}$, and for all $x\in P\setminus P_{\max}$,
  $T_{\lambda(x)}=\bigotimes_{x\ll y} S_{\lambda(y)}$.
\end{definition}
\aain{This suggests that the sources of $\min(P)$ and the targets of $\max(P)$  are empty. They should be arbitrary. UF OK now?}
\aain{OK now}

\ufin{The above requires (event-ordered) parallel product of ipomsets
  and also ordered parallel product of an ordered set of ipomsets.  Both are easy.}

This generalizes the notion of \emph{coherent ST-words}
from~\cite{conf/ramics/AmraneBCFZ24, DBLP:journals/corr/abs-2505-10461}
to partial orders.

\section{ST-sp-pomsets}

\ufin{Blabla serial product $*$ blabla gluing product $\glue$}

A pomset $P$ on an alphabet $\Sigma$ is \emph{series-parallel} (\emph{sp}) if
\begin{itemize}
\item $P=\emptyset$ or $P=a$ for $a\in \Sigma$, or
\item $P=Q\otimes R$ or $P=Q*R$ for sp-pomsets $Q$ and $R$.
\end{itemize}
Every non-empty sp-pomset $P$ is uniquely represented by a term $t(P)$ 
of series and parallel products of singleton pomsets
(up to associativity of $*$ and $\otimes$ and commutativity of $\otimes$).\ufsb{}{citation needed}

We inductively define a translation from an ST-sp-pomset $P$ to a $\Sigma$-labeled ipomset $\flat(P)$
(this generalizes the translation from coherent ST-words to ipomsets
from~\cite{conf/ramics/AmraneBCFZ24, DBLP:journals/corr/abs-2505-10461}):
\begin{itemize}
\item $\flat(\emptyset)=\emptyset$; $\flat(U)=U$ for $U\in \Sigma^\ST$;
\item $\flat(P\otimes Q)=P\otimes Q$; $\flat(P*Q)=P\glue Q$.
\end{itemize}

\ufin{Can this be extended to a translation from \emph{all} ST-pomsets to ipomsets?}

\ufin{Blabla gluing-parallel}

\begin{proposition}
  For any ST-sp-pomset $P$, $\flat(P)$ is gluing-parallel.
\end{proposition}

\begin{proof}
  \ufin{Clear; but should cite the new version of~\cite{arxiv:AikasFJZ22} shortly to appear.}
\end{proof}

Next we inductively define an inverse translation,
from a $\Sigma$-labeled gp-ipomset $P$ to an ST-sp-pomset $\ST(P)$:
\begin{itemize}
\item $\ST(\emptyset)=\emptyset$; 
\item $\ST(a)=a$, \aain{$a$ is not an ST-pomset: $a \notin \Sigma^\ST$} $\ST(\ibullet a)=\ibullet a$, $\ST(a\ibullet)=a\ibullet$,  $\ST(\ibullet a\ibullet)=\ibullet a\ibullet$ for $a\in \Sigma$;
\item $\ST(P\otimes Q)=\ST(P)\otimes \ST(Q)$;
\item $\ST(P\glue Q)=\phi_1(\ST(P))*\phi_2(\ST(Q))$.
\end{itemize}

\ufin{Above, $\phi_1$ and $\phi_2$ are variants of the constructions from (the new version of)~\cite{arxiv:AikasFJZ22}
  which turn interfaces into intervals.  To be written up.}

\begin{proposition}
  \qquad
  For any ST-sp-pomset $P$, $\ST(\flat(P))=P$.
  For any gp-pomset $P$, $\flat(\ST(P))=P$.
\end{proposition}

\begin{proof}
  \ufin{Add proof (easy).}
\end{proof}

\section{Branching ST-automata}

\ufin{Blabla $\sq S$ denotes set of conclists over set $S$.}

\begin{definition}
  A \emph{branching ST-automaton} $(Q, I, F, {\to}, {\branch}, {\merge}, \ev)$
  consists of a finite set $Q$ of states with initial and accepting states $I, F\subseteq Q$,
  sequential transitions ${\to}\subseteq Q\times \Sigma^\ST\times Q$,
  branching and merging transitions ${\branch}\subseteq Q\times \sq Q$, ${\merge}\subseteq \sq Q\times Q$,
  and a state labeling $\ev: Q\to \sq \Sigma$, such that
  \begin{itemize}
  \item for all $q\tto{U} r$, $\ev(q)=S_U$ and $\ev(r)=T_U$;
  \item for all $q\branch [q_1\evord\dotsm\evord q_n]$ and $[q_1\evord\dotsm\evord q_n]\merge q$, $\ev(q)=\bigotimes_{i=1}^n \ev(q_i)$.
  \end{itemize}
\end{definition}

This generalizes both the branching automata of~\cite{DBLP:journals/tcs/LodayaW00}
and the ST-automata of~\cite{conf/ramics/AmraneBCFZ24, DBLP:journals/corr/abs-2505-10461}.

\aain{This does not generalize general branching automata but only fork-acyclic ones:
  those not containing nested forks (and accepting bounded-width sp pomsets).
  In particular, in the second item above $q_i's$ must be different from $q$, otherwise  $\ev(q)=\bigotimes_{i=1}^n \ev(q_i)$ does not hold.}

\ufin{\dots\ except when $\ev(q)=\emptyset$, then we may have $q_i=q$.
  So, if we map branching automata to branching ST-automata with $\ev(q)=\emptyset$ for every $q$ (which is natural),
  then we \emph{do} generalize those.
  
  That being said, there is some thinking to be done about the lack of fork-cyclicity.
  But we could always do $q\tto{\starter{2\ev(q)}{\ev(q)}} q'$ (excuse my abuse of notation)
  and then a fork back to $q\otimes q$, no?}
 \aain{I think, you are right, we should label the states involved in forks and joins by $\emptyset$. The idea is that each branch of a fork is an HDA with start cell having no event active, and one can join when I terminate the events in each branch.
 Is that too restrictive ?
 This will generalise branching automata, you have just to split each sequential transition of a branching autimata into starting and terminating transition in the branching ST-automata.
 Fork-cyclicity is also possible in this case.
}
\ufin{This begs for a generalization of ``automata on graph alphabets''~\cite{DBLP:journals/corr/abs-2602-10036}
  to ``branching automata on monoidal graph alphabets''.}

Languages of branching ST-automata are defined
just as languages of branching automata in~\cite{DBLP:journals/tcs/LodayaW00}.\ufsb{}{needs more detail $\ddot\smile$}
These are, thus, sets of ST-sp-pomsets or, equivalently, sets of gp-ipomsets.

\section{Branching Higher-Dimensional Automata}

\begin{definition}
  A \emph{branching HDA} $(X, \bot, \top, {\branch}, {\merge})$
  consists of a precubical set $X$,
  subsets $\bot, \top\subseteq X$ of initial and accepting cells,
  and branching and merging transitions ${\branch}\subseteq X\times \sq X$, ${\merge}\subseteq \sq X\times X$,
  such that for all $x\branch [x_1\evord\dotsm\evord x_n]$ and $[x_1\evord\dotsm\evord x_n]\merge x$, $\ev(x)=\bigotimes_{i=1}^n \ev(x_i)$.
  Further, for all $x\branch [x_1\evord\dotsm\evord x_n]$, $A\subseteq \ev(x)$, and $\nu\in \{0, 1\}$,
  also $\delta_A^\nu(x)\branch \delta_A^\nu([x_1\evord\dotsm\evord x_n])$,
  and similarly for mergings.
\end{definition}

\ufin{Above, $\delta_A^\nu([x_1\evord\dotsm\evord x_n])$ works by splitting $A$ across the $\ev(x_i)$.}

\ufin{Actually, these subsumption closures of branchings could be added by modifying the base category:
  turning $\sq$ into a ``branching $\sq$'' with coface maps and cobranchings/comergings.
  Still, the notion of track object would change, so this goes beyond~\cite{DBLP:journals/apal/StruthZ26}.}

\ufin{Above, we also allow nullary and unary branchings.
  Nullary branchings are useless, and unary branchings are sequential teleportations.
  (\cite{DBLP:journals/tcs/LodayaW00} require at least two branches.)}

The \emph{operational semantics} of a branching HDA $X=(X, \bot, \top, {\branch}, {\merge})$
is given by the branching ST-automaton $\ST(X)=(X, \bot, \top, {\to}, {\branch}, {\merge}, \ev)$ with
\begin{multline*}
  {\to}=\{(\delta_A^0(x), (\ev(x)\setminus A, \ev(x), \ev(x)), x)\mid A\subseteq \ev(x)\} \\
  \cup \{(x, (\ev(x), \ev(x), \ev(x)\setminus A), \delta_A^1(x))\mid A\subseteq \ev(x)\}.
\end{multline*}
This defines the language of $X$ as the language of $\ST(X)$.
We also give a more direct definition;
below $L_{\textup{nb}}(X)$ denotes the non-branching language of $X$, \ie~ignoring the branching structure.

\begin{definition}
  The \emph{language} of a branching HDA $X$ is $L(X)=\bigcup_{a\in \bot, b\in \top} L_a^b$,
  where $\{L_a^b\mid a, b\in X\}$ are the least fixed point to the equations
  \begin{equation*}
    L_a^b = L_{\textup{nb}}(X)_a^b \cup L_a^b
    \cup \bigcup_{n\ge 0} \bigcup_{\substack{x\branch [x_1\evord\dotsm\evord x_n] \\ [y_1\evord\dotsm\evord y_n]\merge y}}
    L_a^x \glue \big( L_{x_1}^{y_1}\otimes\dotsm\otimes L_{x_n}^{y_n} \big) \glue L_y^b.
  \end{equation*}
\end{definition}

\ufin{So here, $L_a^b$ is the language from $a$ to $b$.
  Is an equivalent definition known for branching/pomset automata?}

\aain{I found a similar notation (denoted $L_{p,q}$ or something similar) several times in \cite{DBLP:conf/fsttcs/Bedon21,DBLP:journals/corr/Bedon15} and in the proof of Theorem 4.12 in \cite{DBLP:journals/tcs/LodayaW00}.
the equations that most closely resemble them are used in Theorem~6 of \cite{DBLP:conf/fsttcs/LodayaW98}.} 

\ufin{There is another definition of the above, likely equivalent:
  split branchings and mergings into pairs $({\branch}, {\merge})$ of the same cardinality,
  and then do the $\bigcup$ only over those pairs.
  Splitting into pairs cannot a priori always be done,
  but any branching HDA can (probably) be turned into one where the pairing is possible.}

\section{Towards Kleene theorem}

\def\ispPos{\mathbf{ispPos}}
\def\iIPos{\mathbf{iIPos}}

\begin{definition}
	The class of \emph{rational gp-languages} is the smallest  subclass of down-closed languages $L\subseteq \ispPos$ that contains all singletons $\{a\}$, $\{\ibullet a\}$, $\{a \ibullet\}$, $\{\ibullet a\ibullet\}$ for $a\in \Sigma$ and is closed with respect to $+$, $*$, $(-)^*$ and $\otimes$.
\end{definition}

\begin{definition}	
	A gp-language is regular if it is a finite approximation of a language of a finite BHDA.
\end{definition}

\begin{definition}
	The class of \emph{rational interval languages} is the smallest  subclass of down-closed languages $L\subseteq \iIPos$ that contains all singletons $\{a\}$, $\{\ibullet a\}$, $\{a \ibullet\}$, $\{\ibullet a\ibullet\}$ for $a\in \Sigma$ and is closed with respect to $+$, $*$, $(-)^*$ and $\otimes$.
\end{definition}

Let $L_i(X)_x^y$ denote the $i$-th approximation of $L(X)_x^y$.

\begin{lemma}
\label{l:hdalang-is-gp}
	The following classes are equal:
	\begin{itemize}
	\item
		The class gp-languages generated from discrete iconclists by $+$, $*$, $(-)^*$.
	\item	
		The class of regular interval languages.
	\item
		The class of rational interval languages.
	\end{itemize}
\end{lemma}
\begin{proof}
	This is shown in \cite{FJSZ_Kleene} + the observation that interval iposets are closed with respect $+$, $*$, $(-)^*$. (So we do not get any gp-posets that are not interval).
\end{proof}

\begin{lemma}
\label{l:pattern}
	Let $Q_n$  be an FSA over alphabet $\Lambda=\{a_i^j\}_{i=\bot,1,\dotsc,n}^{j=1,\dotsc,n,\top}$ with states $\bot,1,\dotsc,n,\top$ and transitions $i\xrightarrow{a_i^j} j$ for $i\in \{\bot,1,\dotsc,n\}$, $i\in \{1,\dotsc,n,\top\}$. Let $e_n$ be a rational expression that describes the language of $Q_n$. Then for every BHDA $X$ having $n$ branchings-mergings:
	\[
		L_{k+1}(X)_x^y=e_n(v_i^j),
	\]
	where
	\begin{itemize}
	\item $v_\bot^i= L_0(X)_x^{s_i}$,
	\item $v_i^j=\bigotimes_\ell L_k(X)_{s_i,\ell}^{t_i,\ell}*L_0(X)_{t_i}^{s_j}$,
	\item $v_i^\top = \bigotimes_\ell L_k(X)_{s_i,\ell}^{t_i,\ell}*L_0(X)_{t_i}^{y}$,
	\item $v_\bot^\top=L_k(X)_x^y$.
	\end{itemize}
\end{lemma}
\begin{proof}
	Routine.
\end{proof}

\begin{lemma}
	For every BHDA $X$, $x,y\in X$ and $k<\infty$, $L_k(X)_x^y$ is gp-rational.
\end{lemma}
\begin{proof}
	This combines Lemmas \ref{l:hdalang-is-gp} and \ref{l:pattern}.
\end{proof}


\begin{proposition}
	Every rational gp-language is regular.
\end{proposition}
\begin{proof}
	This is the hard part of Kleene theorem which becomes easier for BHDA since we can emulate $\epsilon$-transitions by trivial branchings-mergings.
	
	Singleton languages are regular (easy), $+$ is obtained by the disjoint union of BHDA (easy construction). We need to show that gp-regular languages are closed with respect to concatenation, Kleene star and concurrent composition: below.
\end{proof}

\begin{definition}
	Let $X$, $Y$ be BHDA. The parallel composition $X{{}_x^{x'}\parallel_y^{y'}} Y$ is defined as follows.
	\begin{itemize}
	\item As an HDA, $X\parallel Y = X\sqcup Y \sqcup \square^{\ev(x)\otimes\ev(y)} \sqcup \square^{\ev(x')\otimes\ev(y')}$. New top cells are denoted by $s$ and $t$, respectively.
	\item BMs are branchings-mergings of $X$, $Y$, and additionally $(s\branch x \rightrightarrows y\merge t)$
	\end{itemize}
\end{definition} 
\begin{lemma}
	$L_k(X{}_x^{x'}\parallel_y^{y'} Y)_s^t=L_{k-1}(X)_x^{x'}\otimes L_{k-1}(X)_y^{y'}$.
\end{lemma}


\section{Stuff}

\ufin{Also introduce ST-pomset automata and pomset HDAs? following~\cite{DBLP:journals/jlp/KappeBLSZ19}.
  See~\cite{DBLP:conf/fsttcs/Bedon21} for their relation to branching automata.}

\ufin{Conjecture: branching HDAs,
  with the syntactic branching restriction from~\cite{DBLP:journals/jlp/KappeBLSZ19}
  to avoid context-free things,
  are equivalent to usual HDAs with gp-pomset-track objects.}


% \paragraph*{Acknowledgements}



\bibliographystyle{plain}
\bibliography{bibexport}

\end{document}
